\section{Discussion}
\label{sec:discussion}

\subsection{Physiological Mechanics: Channel Attention vs. Spatial Attention in 1D Biosignals}
A primary theoretical insight emerging from our ablation study is the clear divergence between channel attention (SE1D) and spatial attention (CBAM1D) when processing one-dimensional electrophysiological time series. In computer vision, 2D spatial attention effectively isolates regions of interest against redundant background clutter. However, an ECG waveform is fundamentally a continuous, non-stationary temporal sequence where the QRS complex represents an ultra-rapid depolarization spike lasting only $80$--$100$\,ms.

When 1D spatial attention is introduced (as in CBAM1D), temporal pooling and spatial convolutions apply a localized smoothing envelope across the time axis. Consequently, this operation inadvertently suppresses high-frequency steep derivative transitions at the R-peak apex and S-wave nadir, reducing $\Delta\mathrm{SNR}$ down to $9.42$\,dB and downstream S-class F1 to $0.5646$. In contrast, SE1D operates orthogonal to the temporal axis: by re-weighting feature channels, it adaptively amplifies bands encoding cardiac depolarization while damping bands dominated by random electromyographic noise, without modifying the temporal phase or sharpness of individual samples.

\subsection{Clinical Implications of Upstream Mixed-Loss Denoising}
The quantitative NeuroKit2 fiducial measurements (Fig.~\ref{fig:clinical_fiducials}) demonstrate why the loss objective is paramount for clinical translation. Standard MSE loss penalizes amplitude deviations quadratically ($\sim e^2$), prompting the optimizer to compress large peaks toward the mean baseline. In ambulatory monitoring, this peak blunting leads directly to diagnostic errors, such as miscalculating the Sokolow-Lyon index for left ventricular hypertrophy or false-positive alarms for bundle branch blocks.

By integrating absolute error into our mixed objective ($\alpha=0.8$), \HSL{} treats high-amplitude R-peaks with linear penalties, preserving intrinsic voltage levels ($e_R = 0.048$\,mV). In downstream classification, this morphological preservation directly benefits the detection of supraventricular ectopic beats ($\mathbf{S}$). While previous benchmark methods such as Mond{\'e}jar-Guerra \emph{et al.} \cite{mondejar2019} were limited to a positive predictivity of $49.70\%$ on DS2, \RRN{} elevates $+\mathrm{P}$ to \best{72.97\%}. In clinical telemetry units, where alarm fatigue induces severe cognitive strain on nursing staff, an increase of $23.27$ percentage points in precision corresponds to a considerably lower share of false-positive S-class alerts, although the comparison is limited by differences in experimental protocols across studies.

\subsection{Assessment of Hybrid Quantum Computing for Edge Healthcare}
The integration of a 6-qubit parameterized quantum circuit within \RRN{} provides valuable empirical benchmarks for the emerging field of Quantum Machine Learning (QML). Operating with a minimal budget of only \textbf{24 trainable quantum rotation angles}, the hybrid VQC head achieves an ensemble Macro-F1 of \best{0.8176} on the challenging 49{,}298-beat DS2 test set. 

Crucially, the classical residual bypass prevents gradient vanishing associated with barren plateaus \cite{mcclean2018}, ensuring stable convergence during classical backpropagation. While current NISQ quantum hardware does not yet offer quantum supremacy for 1D time-series classification, our results show that an amplitude-encoded VQC with only 24 trainable circuit parameters can approach the accuracy of the classical head ($96.04\%$ vs. $96.43\%$), although the overall head size ($30.9$\,K vs. $37.5$\,K parameters) is comparable because of the classical layers around the circuit.

\subsection{Limitations and Future Trajectories}
Several avenues remain open for future investigation. First, while our noise synthesis utilizes authentic physiological recordings from NSTDB spanning 7 combinations and 10 SNR tiers, real ambulatory environments feature non-stationary artifact transitions. Future work will deploy \HSL{} on uncurated multi-lead datasets (e.g., PTB-XL and CPSC 2018) to test zero-shot cross-database generalization. Second, while the 6-qubit VQC was simulated via state-vector emulators, executing the circuit on physical quantum processors (such as IBM Quantum Heron) with error mitigation techniques (e.g., Zero-Noise Extrapolation) will be evaluated.
